\documentclass[10pt,a4paper]{amsart}
\usepackage[margin=19mm]{geometry}
\usepackage{amsmath,amssymb,mathtools}
\usepackage[scr=rsfs,cal=euler]{mathalfa}
\usepackage{xfrac}
\usepackage[unicode,hidelinks,bookmarks=false]{hyperref}
\newcommand{\ie}{i.e.\ }
\newcommand{\cf}{cf.\ }
\newcommand{\resp}{respectively\ }
\newcommand{\Subsection}{\subsection}
\providecommand{\nobreakdash}{\nobreak\hbox{-}}
\setlength{\emergencystretch}{2em}
\multlinegap0pt
\usepackage{bm}
\usepackage{bbm}
\usepackage{dsfont}
\usepackage{xspace}
\usepackage{cancel}
\usepackage{mathtools}
\usepackage{amsthm}
\makeatletter
\@ifundefined{theorem}{\newtheorem{theorem}{Theorem}}{}
\@ifundefined{lemma}{\newtheorem{lemma}[theorem]{Lemma}}{}
\makeatother
\makeatletter
\expandafter\ifx\csname proof\endcsname\relax
\newenvironment{proof}{\noindent {\bf Proof.}}{\rule{2mm}{2mm}\par\medskip}
\fi
\expandafter\ifx\csname proofof\endcsname\relax
\newenvironment{proofof}{\noindent {\bf Proof of the Theorem 6.1.}}{\rule{2mm}{2mm}\par\medskip}
\fi
\makeatother
\title[Spectral supersaturation: Triangles and bowties]{Spectral supersaturation: Triangles and bowties}
\author{Yongtao Li, Lihua Feng, Yuejian Peng}
\date{Source: arXiv:2407.04950v2 (CC BY 4.0)}
\begin{document}
\maketitle

\noindent\emph{Source and licence.} Spectral supersaturation: Triangles and bowties. Version arXiv:2407.04950v2 (CC BY 4.0), distributed under the Creative Commons Attribution 4.0 International licence, as recorded in the arXiv licence metadata for this exact version and shown on the abstract page https://arxiv.org/abs/2407.04950v2. Full rights notice: licenses/arxiv-2407.04950-CC-BY-4.0.txt.\par\smallskip

\noindent\textit{Editorial scope.} Counted unit: the Lemma whose source label is \texttt{lem-one-edge} in the article (source0.tex lines 1531--1540, source label \texttt{lem-one-edge}), delivered together with the article's own complete original proof of it (source0.tex lines 1542--1609, one \texttt{proof} environment of 68 lines). No mathematics is written, repaired or re-derived: statement and proof are the article's own text line for line. Required context retained uncounted: Lupper (lines 1414--1417); set (lines 1467--1472); lem-W-sub-L (lines 1475--1477). Every definition, display and label of those lines is retained unchanged. Omitted from this package and not redistributed: 1--1413, 1418--1466, 1473--1474, 1478--1530, 1541--1541, 1610--2720. Nothing inside a retained excerpt is omitted or altered: every retained body is delivered exactly as the article's lines are written, so no omission receipt of any kind accompanies this package. Citations: the retained excerpts carry no citation command at all, so no bibliography entry of the article is transcribed and no reference is redistributed. Reference adaptation: none was needed and none was made; every reference inside the delivered unit resolves inside it, so no unresolved reference or citation is produced. Printed slips kept literal: no printed word, symbol, hypothesis or number of the counted statement or of its proof was altered. Layout and class: the article's own class is \texttt{article} with packages including pifont, amsmath, amsfonts, color, float, amssymb, graphicx, bm, enumerate, amsthm, amscd, xy, hyperref; the wrapper is the lane's pinned \texttt{amsart} editorial wrapper at 10pt on A4, which restates the theorem-like environments the retained text needs and adds the article's own macro definitions that the retained text needs (none), with extra wrapper packages (bm, bbm, dsfont, xspace, cancel, mathtools). The delivered numbering is the unit's own. Assets: no retained span contains a figure environment, a graphics-inclusion command, a PDF-page inclusion, a table or a TikZ picture; no image file is copied into this package, the package declares no asset dependency and no image file is redistributed with it. Licence: version arXiv:2407.04950v2 of this article is distributed under the Creative Commons Attribution 4.0 International licence, as recorded in the arXiv licence metadata for this exact version and shown on the abstract page https://arxiv.org/abs/2407.04950v2. A full rights notice is delivered at licenses/arxiv-2407.04950-CC-BY-4.0.txt. Characters: every retained excerpt is pure ASCII, so the delivered engine, XeLaTeX with its default Latin Modern text font, reports no missing character.
\begin{lemma}\label{Lupper}  
We have 
$$|L|< \frac{n}{403}.$$
\end{lemma}

\begin{lemma}\label{set}
Let $A_1, A_2,\dots, A_k$ be $k$  finite sets. Then
\begin{equation*}
\left| \bigcap_{i=1}^k A_i \right| \ge \sum_{i=1}^k |A_i|-(k-1)\left|\bigcup_{i=1}^k A_i\right|.
\end{equation*}
\end{lemma}

  \begin{lemma} \label{lem-W-sub-L}
  We have $W \subseteq L$. 
  \end{lemma}

\begin{lemma}  \label{lem-one-edge}
Both $G[S\setminus L]$ and $G[T\setminus L] $ 
 contain neither three pairwise disjoint edges nor two intersecting edges.   
Furthermore, 
we have $e(S\setminus L) \le 2$ and 
$e(T\setminus L) \le 2$. So 
there exist independent sets
 $I_S\subseteq S\setminus L$  and $I_T\subseteq T\setminus L$ such that
$ |I_S|>  |S|- \frac{n}{402} $ and $|I_T|> |T|- \frac{n}{402} $. 
 \end{lemma}

\begin{proof}
We proceed the proof in two cases. 

{\bf Case 1.}  
If $e_1,e_2$ and $e_3$ are pairwise disjoint edges in $G[S\setminus L]$, 
then we can denote 
$e_1=\{u_1,u_2\}$, $e_2=\{u_3,u_4\}$ 
and $e_3=\{u_5,u_6\}$. 
For $1\le i\le 6 $,   
since $u_i \notin L$, we get  
$d(u_i)> \left(\frac{1}{2}-\frac{1}{200} \right)n$.
By Lemma~\ref{lem-W-sub-L}, 
we have $u_i \notin W$ and 
$d_S(u_i)<  \frac{n}{150}$. Hence
$d_T(u_i)=d(u_i)-d_S(u_i)> \left(\frac{1}{2}-\frac{1}{200}-\frac{1}{150}\right)n.$
By Lemma~\ref{set}, we  get 
    \begin{eqnarray*}
\left |\bigcap_{i=1}^{6}N_T(u_i) \right| &\ge &  \sum_{i=1}^{6}|N_T(u_i)|-5 \left|\bigcup_{i=1}^{6}N_T(u_i) \right| \\
& \ge& \left (\frac{1}{2}-\frac{1}{200}-\frac{1}{150} \right)n\cdot 6 - 
5 \left(\frac{n}{2}+ 14 \sqrt{n} \right) \\
 &>& \frac{n}{6}, 
\end{eqnarray*}
where the last inequality follows by $n\ge 7.1\times 10^{4}$. 
Observe that 
each vertex of the common neighbors of 
$\{u_1,u_2,\ldots ,u_6\}$ leads to three copies of the bowtie. Consequently,  there are 
 more than $\frac{n}{2}$ copies of the bowtie in $G$, 
 which is a contradiction. 

{\bf Case 2.}  
Assume that $e_1$ and $e_2$ are intersecting edges in $G[S\setminus L]$,  
and $e_1,e_2$ form a copy of $K_{1,2}$. 
We denote 
$e_1=\{u_1,u_2\}$ and $e_2=\{u_1,u_3\}$. Similarly, we can see that 
for each $i \in \{2,3\}$, 
   we have
   \begin{eqnarray*}
 |N_T(u_1)\cap N_T(u_i)| 
 &\ge&  |N_T(u_1) | + |N_T(u_i)|  -|T| \\
& \ge & 
2 \left(\frac{1}{2}-\frac{1}{200} -\frac{1}{150} \right)n 
- \left( \frac{n}{2} +14\sqrt{n}\right)\\
 &>&\frac{n}{3}, 
\end{eqnarray*}
where the last inequality holds for $n\ge 9541$. 
In  other words, there are more than 
 $n/3$ choices of a vertex $v_2\in T$ 
such that $u_1u_2 v_2$ forms a triangle. 
Similarly, we have at least $n/3$ choices of a vertex 
$v_3\in T$ such that $v_3\neq v_2$ and 
$u_1u_3v_3$ forms a triangle. 
Thus, we can find $\frac{1}{2}\frac{n}{3} (\frac{n}{3}-1) $ copies of the bowtie in $G$, a contradiction. 
 
 To sum up, we have proved that 
both $G[S\setminus L]$ and $G[T\setminus L]$ 
have at most two non-trivial components, and each of which consists of a single edge. 
Consequently, we get $e(S\setminus L)\le 2$ and $e(T\setminus L)\le 2$. 
Now, by deleting  
 one vertex of each edge in $G[S\setminus L]$, we can obtain 
 a large independent set. 
 So there exists an independent set
 $I_S\subseteq S\setminus L$ such that
  $
  |I_S|\ge |S\setminus L| - 2 
  > |S| - \frac{n}{402}$ by Lemma \ref{Lupper}. 
 The same argument gives that there  is an independent set $I_T\subseteq T\setminus L$ with
$ |I_T|> |T|- \frac{n}{402}$.  
 \end{proof}

\end{document}
